%\usepackage{graphicx}

\begin{titlepage}
	\clearpage
	\newgeometry{left=6.50cm, right=3.25cm, bottom=4.00cm}
	\fontfamily{phv}
	\begin{figure}[t!]		
		\vspace*{(-2.54cm -\voffset-\topmargin-\headheight-\headsep)}
		\includegraphics[width=1.00\paperwidth,pagecenter]{Images/FrontCoverInnerLogo.png}
	\end{figure}
	
	
	\begin{flushleft}
		%	\raggedright
		\newcommand\nbvspace[1][3]{\vspace*{\stretch{#1}}}
		% allow some slack to avoid under/overfull boxes
		\newcommand\nbstretchyspace{\spaceskip0.5em plus 0.25em minus 0.25em}
		% To improve spacing on titlepages
		\newcommand{\nbtitlestretch}{\spaceskip0.6em}
		\pagestyle{empty}
		
		\bfseries
		\nbvspace[3]		
		{\Large
			TÍTULO DO TEMA EM MAIÚSCULAS}
		
		\mdseries		% Texto norma, não bold
		{\Large SUBTÍTULO DO TEMA (SE APLICÁVEL)}
		
		
		
		\nbvspace[2]
		{\large Nome completo do mestrando } \\[0.6em]
		{\small \emph{ Orientador}} \\[0.2em]
		{\normalsize Nome do Orientador} \\[0.6em]
		{\small \emph{ Coorientador}} \\[0.2em]
		{\normalsize Nome do Orientador}\\[2.0em]
		
		
		
		\small Dissertação apresentada\\
		\small ao Instituto Politécnico do Cávado e do Ave\\
		\small para obtenção do Grau de Mestre em Engenharia em Desenvolvimento em Jogos Digitais \\[2.0em]
		
		{\small \emph{ Este trabalho não inclui as críticas e sugestões feitas pelo Júri.}} \\[3.0em]
		
		
		
		{\normalsize mês, ano }
		\nbvspace[1]
		
	\end{flushleft} 
\end{titlepage}